\documentclass[11pt,a4paper]{article}
\usepackage{amsmath,amssymb,amsthm}
\usepackage[margin=2.5cm]{geometry}
\usepackage{hyperref}

\newtheorem{definition}{Definition}[section]
\newtheorem{theorem}[definition]{Theorem}
\newtheorem{proposition}[definition]{Proposition}
\newtheorem{corollary}[definition]{Corollary}
\newtheorem{remark}[definition]{Remark}
\newtheorem{conjecture}[definition]{Conjecture}

\title{Hyperseed Formalization:\\
  Is ChatGPT Real Progress Toward Human-Level AGI?}
\author{ProtomegaTron (automated formalization)}
\date{2026-09-19 (deepened v2)}

\begin{document}
\maketitle

\begin{abstract}
We formalize Ben Goertzel's ``Is ChatGPT Real Progress Toward Human-Level
AGI?'' (2023-01-18) within the Hyperseed ontology, incorporating d-calculus
extensions from Hyperseed v2. ChatGPT's pattern matching is modeled as a
high-curvature fiber shortcut that bypasses intermediate fiber structure,
its missing capabilities as absent fiber dimensions with undefined gradients,
and scaling as volume inflation that preserves topology without adding the
new dimensions AGI requires.
\end{abstract}

\tableofcontents

\section{Fiber shortcut: pattern matching vs understanding}

\begin{definition}[Full fiber traversal]
\emph{Full fiber traversal} for a cognitive task maps input through a
sequence of fiber stages:
\[
  \sigma_{\text{full}}: F_{\text{input}} \xrightarrow{\phi_1}
    F_{\text{percept}} \xrightarrow{\phi_2} F_{\text{concept}}
    \xrightarrow{\phi_3} F_{\text{reason}} \xrightarrow{\phi_4}
    F_{\text{output}}
\]
Each stage $\phi_i$ adds fiber dimensions: perception adds sensory context,
conceptualization adds meaning and relations, reasoning adds logical
structure and implications. The output carries all intermediate fiber
structure.
\end{definition}

\begin{definition}[Fiber shortcut --- claims 1, 2, 7]
A \emph{fiber shortcut} compresses full traversal into a single mapping:
\[
  \sigma_{\text{short}}: F_{\text{input}} \xrightarrow{\psi}
    F_{\text{output}}
\]
where $\psi$ is learned from training examples of the full traversal
$\sigma_{\text{full}}$. The shortcut approximates the input-output
relation without constructing intermediate fiber.
\end{definition}

\begin{proposition}[Shortcut fragility --- claim 8]
\label{prop:fragility}
Within the training distribution $\mathcal{D}_{\text{train}}$, the shortcut
is indistinguishable from full traversal:
\[
  x \in \mathcal{D}_{\text{train}} \implies
  \sigma_{\text{short}}(x) \approx \pi_{\text{output}} \circ
    \sigma_{\text{full}}(x)
\]
Outside the training distribution, the shortcut fails unpredictably:
\[
  x \notin \mathcal{D}_{\text{train}} \implies
  \|\sigma_{\text{short}}(x) - \pi_{\text{output}} \circ
    \sigma_{\text{full}}(x)\| \text{ can be arbitrarily large}
\]
The failure is often \emph{confident} --- the shortcut produces output
with the same fluency as correct output, because the shortcut's quality
metric is fluency (surface fiber), not correctness (deep fiber).
\end{proposition}

\section{Thin and absent fiber dimensions}

\begin{proposition}[Missing components as fiber deficiency --- claims 3, 9]
\label{prop:thin}
ChatGPT's fiber bundle has the following dimensional profile:
\begin{align*}
  \dim(F_{\text{language}}) &\gg 1 \quad \text{(thick --- well-trained)} \\
  \dim(F_{\text{knowledge}}) &\gg 1 \quad \text{(thick --- broad training data)} \\
  \dim(F_{\text{memory}}) &\approx 0 \quad \text{(absent --- no persistence)} \\
  \dim(F_{\text{embodiment}}) &= 0 \quad \text{(absent --- no physical grounding)} \\
  \dim(F_{\text{reasoning}}) &\ll 1 \quad \text{(thin --- surface patterns only)} \\
  \dim(F_{\text{metacognition}}) &= 0 \quad \text{(absent --- no self-reflection)}
\end{align*}
The system has impressive \emph{breadth} (many base points covered) but
insufficient \emph{depth} (key fiber dimensions are thin or absent).
\end{proposition}

\section{Scaling and topology}

\begin{proposition}[Scaling limits --- claims 4, 6]
\label{prop:scaling}
Scaling (increasing parameters $N$) is useful as a component but cannot
reach AGI alone. ChatGPT-like systems contribute a high-quality
language-processing connection but need composition with other cognitive
connections to form a complete AGI architecture.
\end{proposition}

\section{d-calculus formulation (Hyperseed v2)}

\begin{proposition}[Shortcut curvature --- claim 10]
\label{prop:curvature}
The fiber shortcut has high curvature --- the mapping $\psi$ bends
sharply from input to output:
\[
  \|F_{\nabla_{\text{short}}}\| \gg \|F_{\nabla_{\text{full}}}\|
\]
High curvature implies sensitivity to perturbations: small changes in
input can cause large, unpredictable changes in output (hallucinations,
inconsistencies). Full fiber traversal distributes the mapping across
multiple stages, each with moderate curvature, producing smoother
overall behavior:
\[
  \|F_{\nabla_{\text{full}}}\| \leq \sum_i \|F_{\nabla_{\phi_i}}\|
  \ll \|F_{\nabla_{\text{short}}}\|
\]
\end{proposition}

\begin{proposition}[Missing fiber gradient --- claim 11]
\label{prop:gradient}
In fiber dimensions where $\dim(F_d) = 0$ (memory, embodiment,
metacognition), the gradient is undefined:
\[
  \dim(F_d) = 0 \implies \nabla_d F \text{ is undefined}
\]
This is a \emph{qualitative} deficiency, not a quantitative one. Scaling
increases $\dim(F_d)$ for existing dimensions $d$ but cannot create new
dimensions:
\[
  \text{Scale}(F) = \lambda F \implies \{d : \dim(F_d) > 0\}
  \text{ is unchanged}
\]
AGI requires the creation of new fiber dimensions, not the inflation of
existing ones.
\end{proposition}

\begin{proposition}[Component connection --- claim 12]
\label{prop:component}
ChatGPT contributes a specific connection $\nabla_{\text{lang}}$ that
maps between linguistic base space $B_{\text{ling}}$ and concept fiber
$F_{\text{concept}}$:
\[
  \nabla_{\text{lang}}: TB_{\text{ling}} \to \text{End}(F_{\text{concept}})
\]
This connection is high-quality (well-trained, fluent) but must be
composed with other connections for full AGI:
\[
  \nabla_{\text{AGI}} = \nabla_{\text{lang}} \oplus \nabla_{\text{reason}}
    \oplus \nabla_{\text{mem}} \oplus \nabla_{\text{embody}}
    \oplus \nabla_{\text{meta}}
\]
The composition requires interface connections $\nabla_{\text{interface}}^{ij}$
between each pair of connections --- this is cognitive synergy in
connection-theoretic language.
\end{proposition}

\begin{proposition}[Scaling as volume inflation --- claim 13]
\label{prop:volume}
Scaling preserves fiber topology while inflating volume:
\[
  \text{vol}(F_{\text{scaled}}) = \lambda \cdot \text{vol}(F_{\text{base}}),
  \quad \pi_1(F_{\text{scaled}}) \cong \pi_1(F_{\text{base}})
\]
The fundamental group (topology) is unchanged --- no new loops, no new
dimensions, no new connections. AGI requires topological change:
\[
  \pi_1(F_{\text{AGI}}) \not\cong \pi_1(F_{\text{ChatGPT}})
\]
New fiber dimensions (memory, embodiment, metacognition) create new
topological features that scaling cannot produce.
\end{proposition}



%======================================================================
\section{OCO/2 crosswalk (oco/2:2.0.0-alpha.1)}
\emph{Added 2026-10-03. Interpretive mapping onto the Omega Core Ontology
record registry; OCO/2 is an unfrozen alpha candidate.}

\begin{proposition}[Shortcut and axes --- claims 14--17]
An output whose record is only an Execution result can be scored but not checked step by step. If the capability Assessment is $(a_1,\dots,a_m)$ and scaling raises only the $a_i$ already positive, axes with $a_j\approx 0$ stay near zero; closing them needs a component that produces the missing records.
\end{proposition}

\end{document}
